\documentclass{article}

\usepackage[spanish]{babel}
\usepackage{amsmath,amssymb,tikz}
\usepackage{cancel}

\newcommand{\N}{\mathbb{N}}
\newcommand{\Z}{\mathbb{Z}}
\newcommand{\Q}{\mathbb{Q}}
\newcommand{\R}{\mathbb{R}}
\newcommand{\C}{\mathbb{C}}

\title{Tarea 2 Cálculo II}
\author{Benjamín Rivas Beltrán}
\date{\today}

\begin{document}

\maketitle

\section{Suma áreas triángulo}

Dado un triángulo equilátero de lado $a$.

\begin{center}
    \begin{tikzpicture}
        \draw (0,0) -- (4,0) -- (2, {2 * sqrt(3)}) -- cycle;
        \node at (2,-0.35) {$a$};
        \node at (0.65, {sqrt(3) + 0.15}) {$a$};
        \node at (3.35, {sqrt(3) + 0.15}) {$a$};
    \end{tikzpicture}
\end{center}

Su área se puede calcular trazando una altura y dividiendo el triángulo en dos triángulos rectángulos. La altura se puede calcular usando el teorema de Pitágoras:

\begin{center}
    \begin{tikzpicture}
        \draw (0,0) -- (4,0) -- (2, {2 * sqrt(3)}) -- cycle;
        \draw[dashed] (2,0) -- (2, {2 * sqrt(3)});
        \node at (2,-0.35) {$a$};
        \node at (0.65, {sqrt(3) + 0.15}) {$a$};
        \node at (3.35, {sqrt(3) + 0.15}) {$a$};
        \node at (2.5, {sqrt(3) - 0.15}) {$h$};
        \node at (1.2,0.35) {$\frac{a}{2}$};
        \node at (2.8,0.35) {$\frac{a}{2}$};
    \end{tikzpicture}
\end{center}

Usando el teorema de Pitágoras, tenemos:
\begin{align*}
h^2 + \left(\frac{a}{2}\right)^2 &= a^2 \\
h^2 + \frac{a^2}{4} &= a^2 \\
h^2 &= a^2 - \frac{a^2}{4} \\
h^2 &= \frac{3a^2}{4} \\
h &= \frac{\sqrt{3}a}{2}
\end{align*}

Aplicando la fórmula del área de un triángulo, tenemos:
\begin{align*}
A &= \frac{1}{2} \cdot base \cdot altura \\
&= \frac{1}{2} \cdot a \cdot \frac{\sqrt{3}a}{2} \\
&= \frac{\sqrt{3}a^2}{4}
\end{align*}

Dado que el primer triángulo tiene lado $1$ y cada triángulo siguiente tiene como lado la mitad del lado del triángulo anterior, el lado del $n$-ésimo triángulo es $\frac{1}{2^{n-1}}$. Por lo tanto, el área del $n$-ésimo triángulo es:
\begin{align*}
A_n &= \frac{\sqrt{3}}{4} \left(\frac{1}{2^{n-1}}\right)^2 \\
&= \frac{\sqrt{3}}{4} \cdot \frac{1}{2^{2(n-1)}}\\
&= \frac{\sqrt{3}}{4} \cdot \frac{1}{2^{2n-2}}
\end{align*}

Luego, la suma de las áreas de todos los triángulos es:
\begin{align*}
S = \sum_{n=1}^{\infty} A_n &= \sum_{n=1}^{\infty} \frac{\sqrt{3}}{4} \cdot \frac{1}{2^{2n-2}}\\
&= \frac{\sqrt{3}}{4} \sum_{n=1}^{\infty} \frac{1}{2^{2n-2}}\\
&= \frac{\sqrt{3}}{4} \sum_{n=1}^{\infty} \frac{1}{4^{n-1}}\\
&= \frac{\sqrt{3}}{4} \sum_{n=1}^{\infty} \left(\frac{1}{4}\right)^{n-1}\\
\end{align*}

La cual es una serie geométrica con razón $\frac{1}{4}$, por lo que su suma es:
\begin{align*}
S &= \frac{\sqrt{3}}{4} \cdot \frac{1}{1 - \frac{1}{4}}\\
&= \frac{\sqrt{3}}{4} \cdot \frac{1}{\frac{3}{4}}\\
&= \frac{\sqrt{3}}{4} \cdot \frac{4}{3}\\
&= \boxed{\frac{\sqrt{3}}{3}}
\end{align*}

Por lo tanto, la suma de las áreas de todos los triángulos es $\frac{\sqrt{3}}{3}$.

\section{Serie por sumas parciales}

\[
S = \sum_{n=1}^{\infty} \frac{3}{2n^2 + 2n}
\]

Desarrollando el término general, tenemos:
\begin{align*}
\frac{3}{2n^2 + 2n} &= \frac{3}{2n(n+1)}\\
&= \frac{3}{2} \cdot \frac{1}{n(n+1)}\\
&= \frac{3}{2} \cdot \frac{1+n-n}{n(n+1)}\\
&= \frac{3}{2} \left(\frac{1 + n}{n(n+1)} - \frac{n}{n(n+1)}\right)\\
&= \frac{3}{2} \left(\frac{1}{n} - \frac{1}{n+1}\right)
\end{align*}

Reemplazando:
\begin{align*}
S = \frac{3}{2} \sum_{n=1}^{\infty} \left( \frac{1}{n} - \frac{1}{n + 1} \right)
\end{align*}

Sea $a_n = \frac{1}{n}$, entonces:
\[
S = \frac{3}{2} \sum_{n=1}^\infty \left( a_n - a_{n + 1}\right)
\]

Sea $N \in \N$, y:
\[
S_N = \frac{3}{2} \sum_{n=1}^N \left( a_n - a_{n + 1}\right)
\]

Expandiendo la suma, tenemos:
\begin{align*}
S_N &= \frac{3}{2} \left( (a_1 - a_2) + (a_2 - a_3) + \ldots + (a_{N-1} - a_N) + (a_N - a_{N+1}) \right)\\
&= \frac{3}{2} (a_1 - \cancel{a_2} + \cancel{a_2} - \cancel{a_3} + \cancel{a_3} - \ldots - \cancel{a_N} + \cancel{a_N} - a_{N+1})\\
&= \frac{3}{2} (a_1 - a_{N+1})\\
&= \frac{3}{2} \left(1 - \frac{1}{N+1}\right)\\
&= \frac{3}{2} \cdot \frac{N}{N+1}
\end{align*}

Tomando el límite cuando $N$ tiende a infinito, tenemos:
\begin{align*}
S = \lim_{N \to \infty} S_N &= \lim_{N \to \infty} \frac{3}{2} \cdot \frac{N}{N+1}\\
&= \frac{3}{2} \cdot \lim_{N \to \infty} \frac{N}{N+1}\\
&= \frac{3}{2} \cdot \lim_{N \to \infty} \frac{N}{N \left(1 + \frac{1}{N} \right)}\\
&= \frac{3}{2} \cdot \lim_{N \to \infty} \frac{1}{1 + \frac{1}{N}}\\
&= \frac{3}{2} \cdot \frac{1}{1 + 0}\\
&= \boxed{\frac{3}{2}}
\end{align*}

\[
\therefore
\quad
\sum_{n=1}^{\infty} \frac{3}{2n^2 + 2n} = \frac{3}{2}
\]

\section{Convergencia de series}

\subsection{$\sum_{n=1}^{\infty} \frac{(-2)^{2n+1}}{7^n}$}

Reescribiendo el término general, tenemos:
\begin{align*}
\frac{(-2)^{2n+1}}{7^n} &= \frac{(-2)^{2n} \cdot (-2)}{7^n} \\
&= \frac{(-2)^{2n}}{7^n} \cdot (-2) \\
&= \frac{4^n}{7^n} \cdot (-2) \\
&= \left(\frac{4}{7}\right)^n \cdot (-2)
\end{align*}

Reemplazando:
\begin{align*}
\sum_{n=1}^{\infty} \frac{(-2)^{2n+1}}{7^n} &= \sum_{n=1}^{\infty} \left(\frac{4}{7}\right)^n \cdot (-2) \\
&= -2 \sum_{n=1}^{\infty} \left(\frac{4}{7}\right)^n \\
&= -2 \sum_{n=1}^{\infty} \left(\frac{4}{7}\right)^{n-1+1} \\
&= -2 \sum_{n=1}^{\infty} \left(\frac{4}{7}\right)^{n-1} \cdot \frac{4}{7} \\
&= -2 \cdot \frac{4}{7} \sum_{n=1}^{\infty} \left(\frac{4}{7}\right)^{n-1} \\
&= -\frac{8}{7} \sum_{n=1}^{\infty} \left(\frac{4}{7}\right)^{n-1}
\end{align*}

La cual es una serie geométrica con razón $\frac{4}{7} < 1$, por lo tanto, la serie es convergente.

\subsection{$\sum_{n=1}^{\infty} (\frac{n}{n+1})^{-2n}$}

Sea $a_n = (\frac{n}{n+1})^{-2n}$, entonces:
\begin{align*}
a_n &= \left(\frac{n}{n+1}\right)^{-2n} \\
&= \left(\frac{n+1}{n}\right)^{2n} \\
\end{align*}

Se sabe que para que la serie $\sum_{n=1}^{\infty} a_n$ sea convergente, es necesario que $\lim_{n \to \infty} a_n = 0$. Por lo tanto, calculamos el límite:
\begin{align*}
L = \lim_{n \to \infty} a_n &= \lim_{n \to \infty} \left(\frac{n+1}{n}\right)^{2n} \\
&= \lim_{n \to \infty} e^{\ln(\left(\frac{n+1}{n}\right)^{2n})} \\
&= \lim_{n \to \infty} e^{2n \ln(\frac{n+1}{n})} \\
&= e^{\lim_{n \to \infty} 2n \ln(\frac{n+1}{n})} \\
&= e^{\lim_{n \to \infty} 2n \ln(1 + \frac{1}{n})} \\
&= e^{\lim_{n \to \infty} \frac{\ln(1 + \frac{1}{n})}{\frac{1}{2n}}} 
\end{align*}

Aplicando L'Hôpital:
\begin{align*}
L &= e^{\lim_{n \to \infty} \frac{\frac{-1}{n^2(1 + \frac{1}{n})}}{-\frac{1}{2n^2}}} \\
&= e^{\lim_{n \to \infty} \frac{2}{1 + \frac{1}{n}}} \\
&= e^{\frac{2}{1 + 0}} \\
&= e^2
\end{align*}

Dado que $L = e^2 \neq 0$, la serie $\sum_{n=1}^{\infty} (\frac{n}{n+1})^{-2n}$ es divergente.

\subsection{$\sum_{n=3}^{\infty} n^3 e^{-n^4}$}

Sea $f : [3, \infty)  \to (0, \infty)$ una función dada por:
\begin{align*}
f(x) = x^3 e^{-x^4}
\end{align*}

Note que, para $x \geq 3$, se tiene que $f(x) > 0$. Pues $e^{-x^4} > 0$ para todo $x \in \R$, y $x^3 > 0$ para todo $x > 0$. Además, note que $f$ es una función continua, pues es el producto de funciones continuas.

Derivando $f$, tenemos:
\begin{align*}
f'(x) &= 3x^2 e^{-x^4} + x^3 \cdot (-4x^3) e^{-x^4} \\
&= 3x^2 e^{-x^4} - 4x^6 e^{-x^4} \\
&= e^{-x^4} (3x^2 - 4x^6) \\
&= x^2 e^{-x^4} (3 - 4x^4)
\end{align*}

Note que, para $x \geq 3$, se tiene que:
\begin{align*}
x^4 &\geq 81 \\
4x^4 &\geq 324 \\
-4x^4 &\leq -324 \\
3 - 4x^4 &\leq 3 - 324 \\
3 - 4x^4 &\leq -321 < 0
\end{align*}

Por lo tanto, $f'(x) < 0$ para todo $x \geq 3$, lo que implica que $f$ es una función decreciente en el intervalo $[3, \infty)$.

Dado que $f$ es una función positiva, continua y decreciente en el intervalo $[3, \infty)$, se puede aplicar el criterio de la integral para determinar la convergencia de la serie $\sum_{n=3}^{\infty} n^3 e^{-n^4}$.
Calculando la integral impropia, tenemos:
\begin{align*}
\int_3^{\infty} x^3 e^{-x^4} dx &= \lim_{t \to \infty} \int_3^t x^3 e^{-x^4} dx
\end{align*}

Haciendo el cambio de variable $u = -x^4$, tenemos $du = -4x^3 dx$, lo que implica que $-du/4 = x^3 dx$. Además, cuando $x = 3$, se tiene que $u = -81$, y cuando $x = t$, se tiene que $u = -t^4$. Por lo tanto:
\begin{align*}
\int_3^t x^3 e^{-x^4} dx &= \int_{-81}^{-t^4} e^u \cdot \left(-\frac{du}{4}\right) \\
&= -\frac{1}{4} \int_{-81}^{-t^4} e^u du \\
&= -\frac{1}{4} \left( e^{-t^4} - e^{-81} \right) \\
&= \frac{1}{4} \left( e^{-81} - e^{-t^4} \right)
\end{align*}

Tomando el límite cuando $t$ tiende a infinito, tenemos:
\begin{align*}
\int_3^{\infty} x^3 e^{-x^4} dx &= \lim_{t \to \infty} \frac{1}{4} \left( e^{-81} - e^{-t^4} \right) \\
&= \frac{1}{4} \left( e^{-81} - 0 \right) \\
&= \frac{e^{-81}}{4} \\
&= \boxed{\frac{1}{4 \cdot e^{81}}}
\end{align*}

Dado que la integral impropia $\int_3^{\infty} x^3 e^{-x^4} dx$ es convergente, se concluye que la serie $\sum_{n=3}^{\infty} n^3 e^{-n^4}$ también es convergente.

\section{Intervalo de convergencia}

\begin{align*}
\sum_{n=1}^{\infty} \frac{(-1)^{n+1} \cdot x^n}{2^n \cdot (n+3)}
\end{align*}

Sea $x \in \R$, se puede aplicar el criterio de la raíz para determinar el intervalo de convergencia de la serie. Calculando el límite, tenemos:
\begin{align*}
L = \lim_{n \to \infty} \sqrt[n]{\left|\frac{(-1)^{n+1} \cdot x^n}{2^n \cdot (n+3)}\right|} &= \lim_{n \to \infty} \sqrt[n]{\frac{|x|^n}{2^n \cdot (n+3)}} \\
&= \lim_{n \to \infty} \frac{|x|}{2} \cdot \sqrt[n]{\frac{1}{n+3}} \\
&= \frac{|x|}{2} \cdot \lim_{n \to \infty} \sqrt[n]{\frac{1}{n+3}} \\
&= \frac{|x|}{2} \cdot \lim_{n \to \infty} \frac{1}{\sqrt[n]{n+3}} \\
&= \frac{|x|}{2} \cdot \frac{1}{\lim_{n \to \infty} \sqrt[n]{n+3}} \\
&= \frac{|x|}{2} \cdot \frac{1}{\lim_{n \to \infty} e^{\ln \sqrt[n]{n+3}}} \\
&= \frac{|x|}{2} \cdot \frac{1}{\lim_{n \to \infty} e^{\frac{1}{n} \ln (n+3)}} \\
&= \frac{|x|}{2} \cdot \frac{1}{e^{\lim_{n \to \infty} \frac{1}{n} \ln (n+3)}} \\
&= \frac{|x|}{2} \cdot \frac{1}{e^{\lim_{n \to \infty} \frac{\ln (n+3)}{n}}} \\
&= \frac{|x|}{2} \cdot \frac{1}{e^0} \\
&= \frac{|x|}{2}
\end{align*}

Dado que $L = \frac{|x|}{2}$, la serie es convergente si $L < 1$, lo que implica que $|x| < 2$. Por lo tanto, la serie es convergente para $x \in (-2, 2)$.

Para determinar la convergencia en los extremos del intervalo, se evalúa la serie en $x = -2$ y $x = 2$.

Cuando $x = -2$, el término general de la serie es:
\begin{align*}
\frac{(-1)^{n+1} \cdot (-2)^n}{2^n \cdot (n+3)} &= \frac{(-1)^{n+1} \cdot (-1)^n \cdot 2^n}{2^n \cdot (n+3)} \\
&= \frac{(-1)^{ 2n+1} \cdot 2^n}{2^n \cdot (n+3)} \\
&= \frac{-1}{n+3}
\end{align*}

La serie resultante es $\sum_{n=1}^{\infty} \frac{-1}{n+3}$, la cual es divergente, pues es una serie armónica desplazada.

Cuando $x = 2$, el término general de la serie es:
\begin{align*}
\frac{(-1)^{n+1} \cdot 2^n}{2^n \cdot (n+3)} &= \frac{(-1)^{n+1}}{n+3}
\end{align*}

La serie resultante es $\sum_{n=1}^{\infty} \frac{(-1)^{n+1}}{n+3}$, la cual es convergente, pues es una serie alternante con términos que tienden a cero, por lo que se cumple el criterio de Leibniz.

Por lo tanto, el intervalo de convergencia de la serie es $(-2, 2]$.

\section{Expansión de $\ln(1-x^2)$}

Sea $f : (-1, 1) \to \R_{\leq 0}$ una función dada por:
\begin{align*}
f(x) = \ln(1-x^2)
\end{align*}

Factorizando y aplicando la propiedad de los logaritmos, tenemos:
\begin{align*}
f(x) &= \ln((1-x)(1+x)) \\
&= \ln(1-x) + \ln(1+x)
\end{align*}

Derivando $f$, tenemos:
\begin{align*}
f'(x) &= \frac{1}{1-x} \cdot (-1) + \frac{1}{1+x} \cdot (1) \\
&= -\frac{1}{1-x} + \frac{1}{1+x} \\
&= \frac{1}{1+x} - \frac{1}{1-x} \\
&= \frac{1}{1-(-x)} - \frac{1}{1-x} \\
\end{align*}

Note que, para $x \in (-1, 1)$, se tiene que $-x \in (-1, 1)$. Por lo tanto, se puede aplicar la expansión de serie geométrica para cada término:

\begin{align*}
f'(x) &= \sum_{n=0}^{\infty} (-x)^n - \sum_{n=0}^{\infty} x^n \\
&= \sum_{n=0}^{\infty} (-1)^n x^n - \sum_{n=0}^{\infty} x^n \\
&= \sum_{n=0}^{\infty} \left( (-1)^n - 1 \right) x^n
\end{align*}

Note que, para $n$ par, se tiene que $(-1)^n - 1 = 0$, y para $n$ impar, se tiene que $(-1)^n - 1 = -2$. Por lo tanto, se pueden tomar solo los términos impares de la serie:
\begin{align*}
f'(x) &= \sum_{n=0}^{\infty} \left((-1)^{2n+1} - 1\right) x^{2n+1} \\
&= \sum_{n=0}^{\infty} (-2) x^{2n+1} \\
&= -2 \sum_{n=0}^{\infty} x^{2n+1} \\
\end{align*}

Integrando ambos lados, tenemos:
\begin{align*}
\int f'(x) dx &= \int -2 \sum_{n=0}^{\infty} x^{2n+1} dx \\
f(x) &= -2 \sum_{n=0}^{\infty} \int x^{2n+1} dx \\
&= -2 \sum_{n=0}^{\infty} \frac{x^{2n+2}}{2n+2} + C \\
&= - \sum_{n=0}^{\infty} \frac{x^{2n+2}}{n+1} + C\\
&= - \sum_{n=1}^{\infty} \frac{x^{2n}}{n} + C
\end{align*}

Dado que $f(0) = \ln(1-0^2) = 0$, se tiene que:
\begin{align*}
0 &= f(0) \\
0 &= - \sum_{n=1}^{\infty} \frac{0^{2n}}{n} + C \\
0 &= - \sum_{n=1}^{\infty} 0 + C \\
0 &= C
\end{align*}

Por lo tanto, la expansión de $\ln(1-x^2)$ es:
\begin{align*}
\ln(1-x^2) &= - \sum_{n=1}^{\infty} \frac{x^{2n}}{n} \\
&= -x^2 - \frac{x^4}{2} - \frac{x^6}{3} - \frac{x^8}{4} - \ldots
\end{align*}

\section{Serie de Taylor de $\sin(2x)$ alrededor de $x = \frac{\pi}{2}$}

Sea $f : \R \to \R$ una función dada por:
\begin{align*}
f(x) = \sin(2x)
\end{align*}

Derivando $f$, tenemos:
\begin{align*}
f'(x) &= 2 \cos(2x) \\
f''(x) &= -4 \sin(2x) \\
f'''(x) &= -8 \cos(2x) \\
f^{(4)}(x) &= 16 \sin(2x) \\
f^{(5)}(x) &= 32 \cos(2x) \\
f^{(6)}(x) &= -64 \sin(2x) \\
f^{(7)}(x) &= -128 \cos(2x) \\
f^{(8)}(x) &= 256 \sin(2x) \\
\end{align*}

Las derivadas de $f$ siguen un patrón si separan los casos en que $n$ es par o impar. Para $n$ par, se tiene que:
\begin{align*}
f^{(n)}(x) &= (-1)^{\frac{n}{2}} 2^n \sin(2x)
\end{align*}

Y para $n$ impar, se tiene que:
\begin{align*}
f^{(n)}(x) &= (-1)^{\frac{n-1}{2}} 2^n \cos(2x)
\end{align*}

Reemplazando $x = \frac{\pi}{2}$, tenemos:
\begin{align*}
f^{(n)}\left(\frac{\pi}{2}\right) &= \begin{cases}
(-1)^{\frac{n}{2}} 2^n \sin(\pi) & \text{si } n \text{ es par} \\
(-1)^{\frac{n-1}{2}} 2^n \cos(\pi) & \text{si } n \text{ es impar}
\end{cases} \\
&= \begin{cases}
0 & \text{si } n \text{ es par} \\
(-1)^{\frac{n-1}{2}} 2^n \cdot (-1) & \text{si } n \text{ es impar}
\end{cases} \\
&= \begin{cases}
0 & \text{si } n \text{ es par} \\
(-1)^{\frac{n+1}{2}} 2^n & \text{si } n \text{ es impar}
\end{cases} 
\end{align*}

Dado que $f^{(n)}\left(\frac{\pi}{2}\right) = 0$ para todo $n$ par, solo se tomarán los términos impares de la serie de Taylor. Por lo tanto, la serie de Taylor de $\sin(2x)$ alrededor de $x = \frac{\pi}{2}$ es:
\begin{align*}
\sum_{n=0}^{\infty} \frac{f^{(2n+1)}\left(\frac{\pi}{2}\right)}{(2n+1)!} \left(x - \frac{\pi}{2}\right)^{2n+1} &= \sum_{n=0}^{\infty} \frac{(-1)^{n+1} 2^{2n+1}}{(2n+1)!} \left(x - \frac{\pi}{2}\right)^{2n+1} \\
&= -2 \sum_{n=0}^{\infty} \frac{(-1)^n 2^{2n}}{(2n+1)!} \left(x - \frac{\pi}{2}\right)^{2n+1} \\
&= -2 \sum_{n=0}^{\infty} \frac{(-1)^n 4^n}{(2n+1)!} \left(x - \frac{\pi}{2}\right)^{2n+1}
\end{align*}

Finalmente, la serie de Taylor de $\sin(2x)$ alrededor de $x = \frac{\pi}{2}$ es:
\begin{align*}
\sin(2x) &= -2 \sum_{n=0}^{\infty} \frac{(-1)^n 4^n}{(2n+1)!} \left(x - \frac{\pi}{2}\right)^{2n+1} \\
&= -2 \left( \frac{1}{1!} \left(x - \frac{\pi}{2}\right) - \frac{4}{3!} \left(x - \frac{\pi}{2}\right)^3 + \frac{16}{5!} \left(x - \frac{\pi}{2}\right)^5 \mp \ldots \right)
\end{align*}

\end{document}

